\section{Threats to Validity}
\label{sec:threats}

\paragraph{Internal.} The main FlakeAgent runs are single-seeded; agent
decisions are stochastic, so single-run rates carry sampling noise, which we
mitigate with a 3-seed variance subset (Section~\ref{sec:setup}) and fixed
temperature and prompts. The executor is large, hand-built infrastructure;
bugs in strategy implementations or log parsing could bias results, which we
mitigate by grounding strategies against the dataset's own mechanisms (T2
reuses iDFlakies, T12 reuses NonDex) and releasing all logs.

\paragraph{External.} The benchmark is drawn from ReproFlake, which is
Java/Maven-only and skewed toward implementation-dependent tests in its full
distribution; our subset deliberately balances categories, so aggregate
numbers are not a statement about ReproFlake's natural distribution. All 116
tests are within one dataset, so results may not transfer to other ecosystems
or proprietary code. All LLM calls use vendor-hosted 2026 models; endpoint
behavior may drift after our runs, which we pin in the artifact.

\paragraph{Construct.} ``Reproduction'' is operationalized as a target-test
failure within budget, judged by the observation probe; a probe parsing bug
could miscount. A subtler threat: fault-injection strategies can make a test
fail for reasons \emph{orthogonal} to its actual flakiness — e.g., NonDex
(T12) can surface a reflection-linkage error in an order-dependent test, and
an aggressive sleep (T6) or heap cap (T9) can trip an unrelated timeout. Such
a ``failure'' counts as a reproduction even though it is not \emph{the} flaky
failure; we therefore flag, in the artifact, every reproduction whose
diagnosed category disagrees with its ground-truth label. Diagnosis grounding
is measured against the developer fix diff, which may inflate or deflate
precision, bounded by our additional manual assessment. Where we compare
against baselines we use McNemar's test; model comparison is descriptive
(one run per model), and costs are at provider list prices.
